\section{Multi-lane verified retrieval}
\label{sec:lanes}

All answering agents reach the platform through one chat endpoint. The
server-side \emph{lane policy} inspects the request (pinned libraries,
prompt suffixes, engine) and attaches the corresponding tool whitelist and
instructions; a librarian-lane agent, for example, is denied the vector and
two-stage tools at the protocol level, so the lane is a property of the
enforcement layer, not of prompt discipline. Fig.~\ref{fig:lanes} summarizes
the three lanes.

\begin{figure}[t]
\centering
\includegraphics[width=\textwidth]{fig2_lanes}
\caption{The three retrieval lanes. Every lane ends at the same
content-verification gate and the same honesty contract: zero verified
evidence escalates (Lane A$\to$B) or terminates in an explicit not-found
report.}
\label{fig:lanes}
\end{figure}

\subsection{Why three lanes --- and when to use which}
\label{sec:lanechoice}
The three lanes are not three implementations of one idea; they occupy three
points on a latency/coverage frontier (Table~\ref{tab:lanes}). The Hybrid lane
is the fastest and the most auditable --- it moves orchestration server-side
and ships its audit trail in one response --- and is the interactive default.
The Librarian lane is the one that must not miss a document: it spends its
twelve-call budget on catalog navigation and buys exhaustive coverage with
explicit blind-spot disclosure. The Search lane sits between them as the
latency-tuned interactive path for pinned libraries.

\begin{table}[!htbp]
\centering
\scriptsize
\setlength{\tabcolsep}{3.5pt}
\caption{The three lanes as operating points. Measured figures are the
round-1 agentic runs of Section~\ref{sec:agentic}. All three share one
verification gate and one honesty contract; they differ only in retrieval
affordance and budget.}
\label{tab:lanes}
\begin{tabular}{p{1.85cm}p{3.05cm}p{2.15cm}p{2.25cm}p{3.15cm}}
\toprule
Lane & Optimises for & Tool budget & Measured (round 1) & Use when \\
\midrule
A. Search & latency on pinned libraries & 20-turn cap & 6--8 turns; 70--165\,s & interactive default; single-library questions \\
B. Librarian & exhaustive recall across all libraries & hard 12-call budget & 5--6 turns; 54--91\,s & must-not-miss; enumeration questions; no library pinned \\
C. Hybrid & single-call server-side orchestration & exactly 1 call & 3--5 turns; 36--70\,s & fastest grounded answer; latency-critical \\
\bottomrule
\end{tabular}
\end{table}

Two honest caveats accompany this table. First, in the measurements of
Section~\ref{sec:agentic} the Search lane does not occupy a \emph{distinct}
operating point --- the Hybrid lane is faster and the Librarian lane is more
exhaustive --- so a two-lane deployment (Hybrid plus Librarian) is defensible
today, with the Search lane retained as the pinned-library interactive
default; query-shape auto-selection is future work
(Section~\ref{sec:threats}). Second, and more importantly, the three lanes
differ only in retrieval affordance and budget: they share one gate, one
fail-closed contract, and one abstention semantics. The trust model is a
property of the platform, not of the lane.

\subsection{Lane A: Search --- vector fast lane with a verification gate}
Lane A optimizes for latency on pinned libraries. Phase A0 rewrites the
query into (subject, attribute, constraints) form, splits multi-entity and
comparison questions into parallel sub-queries, and routes incident-style
queries to the experience store (Section~\ref{sec:system}) first. Phase A1
issues one wide dense query
(top-30, threshold 0.35, cross-library balancing), deduplicates by document,
and \emph{keeps all candidate documents rather than truncating to the top
chunks}. Phase A2 submits citation-only references (library id + document
path) to the verification gate in batches of at most six; the server fetches the
full text itself, so the agent never fabricates content when arguing for a
document. A document with vector score 0.95 that the judge rejects is
dropped; a document at 0.40 with one \emph{yes} verdict survives. Phase A3
returns the judge's evidence pack (the only full-content channel in this
lane) and the agent composes an answer whose citations carry confidence
levels (P0 direct evidence, P1 partial, P2 weak).

\subsection{Lane B: Librarian --- hierarchical navigation under a budget}
Lane B is for exhaustive recall across all libraries and for queries whose
vocabulary may not match the corpus. It walks a fixed spine with a hard
budget of twelve tool calls: L0 reads the library catalog (ids, names,
descriptions, counts) and prunes shelves (libraries) by the rewritten
query; L1 lists
each surviving shelf's documents with descriptions; L2 routes by query
shape --- literal/enumeration questions (``which documents mention X'') go
to a machine-side scanner ($\approx$0.5\,s per document, no LLM turns),
semantic questions go to judging; L3 judges in batches of at most six
references (larger batches historically triggered transport timeouts and
fail-closed batch loss); L4 reads only survivors (bounded to 2\,000
characters, starting at the judge-reported evidence line); L5 emits a
five-part report (search paths, answer, sources, confidence, blind spots).
Two rules are absolute: catalog descriptions never prune documents by
themselves, and an exhausted budget terminates in blind-spot disclosure,
not in a guessed answer.

\subsection{Lane C: Hybrid --- server-side orchestration in one call}
Agentic latency is dominated by LLM turns, not by tools (Section
\ref{sec:composition}); Lane C spends exactly one call by moving the entire
pipeline into the server. A single \texttt{kb\_hybrid\_search} call executes
P1 (a parallel block: vector channel, top-30 at threshold 0.35, and catalog
keyword channel); P3 merges candidates by (library, document) with
agreed-by-both
documents first, capped at twelve; P4 judges every candidate with the
verification gate in batches of six (criterion \emph{evidence}, threshold
0.5); P4b
applies a document-level relative cut --- keep scores $\geq \mathit{best} -
0.15$ with a top-5 floor, plus a cross-domain anchor guard that prevents an
inflated out-of-domain score from setting the cut line; and P5 reads the
top five documents in full ($\leq$2\,000 characters each). The response
carries the complete audit trail (per-lane candidates, merge counts, judge
statistics, cut decisions, read contents), so a one-call lane is still an
inspectable lane.

\subsection{The verification gate}
\label{sec:gate}
The gate is a locally hosted evidence-classification model (the Laya
judging service \cite{laya}) invoked through a
\texttt{kb\_laya\_judge} tool. Its contract: given the
query and a batch of document references, fetch each document server-side,
segment it, and score \emph{does this text contain direct evidence that
answers the query} under one of three criteria --- \emph{evidence}
(specific substantiation; mere topic match is insufficient),
\emph{instance} (contains at least one member of the requested class;
routed automatically for enumeration questions), or \emph{auto}. Scores
live in $[0,1]$ with an acceptance threshold of 0.5 (0.3 for enumeration).
All operating constants in this section --- thresholds, the six-reference
judge batch, the hybrid candidate cap of twelve, and the 0.15 relative-cut
margin with its top-5 floor --- are the platform's shipped production
defaults, fixed during earlier deployments on a different corpus, and were
not tuned on the evaluation questions of Section~\ref{sec:setup}. At the
protocol layer the platform documents scores as a 0--8 rubric (topic 0--3,
scenario 0--3, evidence 0--2; $\geq$6 answers directly, 5 escalates to the
librarian lane, $\leq$4 drops); the $[0,1]$ score is the machine-readable
projection of that rubric, and the rubric is the human-auditable contract
published to platform clients. The projection can saturate --- the logged
maximum score is 1.0000 --- one more reason thresholds carry no
discriminating power between golds and same-domain confusers. The judge
model itself is a locally
fine-tuned evidence classifier (checkpoint reference released with the
platform artifacts); it is deterministic in our replays --- repeat calls
with identical inputs reproduced identical orderings
(Section~\ref{sec:retrieval}) --- and it is a \emph{single} checkpoint:
every verifier finding in this paper inherits that choice, which
Section~\ref{sec:threats} treats as an explicit scope limit.
Fail-closed semantics apply to every
failure mode: fetch errors, engine errors, and timeouts count as
rejection. Survivors are aggregated into an evidence pack with
per-passage provenance.

\subsection{Honest abstention}
If a lane ends with zero verified documents, Lane A escalates once to Lane
B; if Lane B also exhausts its budget without a survivor, the answer is the
not-found report: what was searched (libraries, document counts), what the
query was routed as, and which topics sit outside the corpus. Because
abstention is a terminal state in the lane contract --- enforced by the
server-side policy rather than advised in a prompt --- it can be evaluated
like any other outcome, which Section~\ref{sec:agentic} does.

